\documentclass[12pt,a4paper]{report}
\usepackage[utf8]{inputenc}
\usepackage{amsmath}
\usepackage{amsfonts}
\usepackage{amssymb}
\usepackage{amsthm}
\usepackage{hyperref}

\usepackage{multicol}
\usepackage{fancyhdr}
\usepackage[inline]{enumitem}
\usepackage{tikz}
\usepackage{tikz-cd}
\usetikzlibrary{calc}
\usetikzlibrary{shapes.geometric}
\usepackage[margin=0.5in]{geometry}
\usepackage{xcolor}
\usepackage{ esint }

\hypersetup{
    colorlinks=true,
    linkcolor=blue,
    filecolor=magenta,      
    urlcolor=cyan,
    pdftitle={Functional Analysis},
    pdfpagemode=FullScreen,
    }

%\urlstyle{same}

\newcommand{\CLASSNAME}{Comprehensive Exam -- Real Analysis}
\newcommand{\STUDENTNAME}{Paul Carmody}
\newcommand{\ASSIGNMENT}{NOTES }
\newcommand{\DUEDATE}{April 2026}
\newcommand{\SEMESTER}{Spring 2026}
\newcommand{\SCHEDULE}{TBD}
\newcommand{\ROOM}{Crawford 332}

\newcommand{\MMN}{M_{m\times n}}
\newcommand{\FF}{\mathcal{F}}

\pagestyle{fancy}
\fancyhf{} \chead{ \fancyplain{}{\CLASSNAME} }
%\chead{ \fancyplain{}{\STUDENTNAME} }
\rhead{\thepage}
\fancyfoot[C]{-- \thepage --}
\input{../../MyMacros}

\newcommand{\RED}[1]{\textcolor{red}{#1}}
\newcommand{\BLUE}[1]{\textcolor{blue}{#1}}
\newcommand{\FFF}{\mathcal{F}}

\begin{document}

\begin{center}
	\Large{\CLASSNAME -- \SEMESTER} 
\end{center}
\begin{center}
	\STUDENTNAME \\
	\ASSIGNMENT -- \DUEDATE\\
\end{center} 

\begin{enumerate}

\item Let $\mathcal{F}$ be a colletion of measurable functions $f: X \to [0,\infty]$,
\begin{align*}
	c = \inf_{f \in \mathcal{F}} \int f.
\end{align*}Assume that every sequence $\{f_n\}\subset \mathcal{F}$ has a subsequence $f_n \to f$ for some $f \in \mathcal{F}$.  Show that there is an $f \in \mathcal{F}$ with $\int f =c$.

\item Let $f \in L^1, E_1\subset E_2\subset \cdots$ be measurable with $X = \bigcup_n E^n$.  Shown that $\int_{E_n} f \to \int_X f$.

\item If $p \in (r,s) \subset [1,\infty], f \in L^r\cap L^s$, then $f \in L^p$ and $|f|_p \le \max\{|f|_r, |f|_s\}$.

\item If $p \in [1,\infty), f_n \to f$ in $L^p(X), g_n\to g$ a.e., on $X$ and $g_n$ is bounded in $L^\infty(X)$, then $f_ng_n \to fg$ in $L^p(X)$.

\item If $p,q \in (1, \infty), f_n \to f$ in $L^p, g_n\to g$ in $L^q$ where $p^{-1}+q^{-1}=1$, then $f_n g_n \to fg $.
\end{enumerate}

\end{document}
